\documentclass[xcolor=svgnames]{beamer}
\usepackage{latexsym,amssymb,amsmath,stmaryrd}
\usepackage[utf8]{inputenc}
\usepackage[T1]{fontenc}
\usepackage[francais]{babel}
\usepackage{pgf}
\usepackage{tikz}
\usetikzlibrary{arrows,patterns,plotmarks,shapes,snakes,er,3d,automata,%
backgrounds,topaths,trees,petri,mindmap}
\usepackage{pgfbaseimage}
\usepackage{ifpdf}
\ifpdf
\usepackage{graphicx}
\else
%\usepackage[dvips]{graphicx}
\usepackage{pstricks,pst-tree,pst-node}
\newpsobject{showgrid}{psgrid}{subgriddiv=1,griddots=10,gridlabels=6pt}
\fi
\usepackage{moreverb}
\usepackage[normalem]{ulem}
\usepackage{version}
\usepackage{url}
\usepackage{multicol}
\usepackage{listings}
\lstset{
  language=[ANSI]C, 
%  gobble=2, 
%  escapeinside="", 
  basicstyle=\ttfamily, 
%  directivestyle=\color{Fuchsia},
  identifierstyle = \color{DarkOrange},
  keywordstyle=\color{DarkGreen}, 
  commentstyle=\color{red}, 
%  numbers=left, 
%  numbersep=5pt, 
%  numberstyle=\scriptsize, 
%  moredelim=[is][\color{gray}\itshape]{/*}{*/}, 
%  moredelim=[is][\alert]{/+}{+/}, 
%  morecomment=[is]{/=}{=/}, 
%  morecomment=[in]{comment=}{,} 
%  fancyvrb=true 
} 

\newcommand{\binaire}[1]{\ensuremath{\underline{#1}}}
\newcommand{\C}[1]{\texttt{#1}}
\newcommand{\bbbn}{\ensuremath{\mathbb{N}}}
\input{communs/cc_beamer}
\usetheme{classic}
\newcommand{\nowrite}{\put(10,-4){\includegraphics[scale=.05]{creative_commons/nopencil}}}
\newcommand{\youwrite}{\put(10,-4){\includegraphics[scale=.05]{creative_commons/pencil}}}
\newcommand{\aemporter}{\put(10,-6){\includegraphics[scale=.05]{creative_commons/szymonraj_Shopping_bag}}}


%%% Titre -- Cours 1
\title{Éléments d'informatique -- Cours 2. Éléments de systèmes
  d'exploitation. Premier programme C}
\author{Pierre Boudes}
\date{\today}

\begin{document}

%% Page de titre et licence CC.
\begin{frame}
	\titlepage
	\vfill
	\begin{center}
		\CcGroupByNcSa{0.83}{0.95ex}\\[2.5ex]
		{\tiny\CcNote{\CcLongnameByNcSa}}
		\vspace*{-2.5ex}
	\end{center}
\end{frame}

  \section*{Contenu du cours}
\begin{frame}
  \frametitle{Contenu du cours (rappel)}
  \begin{itemize}
  \item Éléments d'architecture des ordinateurs (+mini-assembleur)
  \item \alert{Éléments de systèmes d'exploitation}
\item Programmation structurée impérative (éléments de langage C)
\begin{itemize}
  \item Structure d'un programme C
  \item Variables : déclaration (et initialisation), affectaction
  \item Évaluation d'expressions
  \item Instructions de contrôle : if, for, while
  \item Types de données : entiers, caractères,
    réels, tableaux, enregistrements
  \item Fonctions d'entrées/sorties (scanf/printf)
  \item Écriture et appel de fonctions
  \item Débogage
\end{itemize}
\item Notions de compilation
  \begin{itemize}
  \item Analyse lexicale, analyse syntaxique, analyse sémantique
  \item préprocesseur du compilateur C (include, define)
   \item Édition de lien
  \end{itemize}
\item Algorithmes élémentaires
\item Méthodologie de résolution, manipulation sous linux
\end{itemize}
\end{frame}

\section[Plan]{}
\frame[label=plan]{\tableofcontents%[pausesections]
}

\section{Le(s) système(s) d'exploitation}
\begin{frame}
\frametitle{Le(s) système(s) d'exploitation}

Il n'y a pas de définition précise.\pause
\begin{definition}[système d'exploitation]
  Un \alert{système d'exploitation} est un ensemble de programmes qui
  servent d'interface entre le matériel et les logiciels applicatifs. 
\end{definition}
\pause
Nous distinguerons trois rôles endossés par le SE :
\pause
  \begin{itemize}
  \item Pilotage et abstraction du matériel\pause
  \item Gestionnaire et ordonnanceur de processus\pause
  \item Environnement de travail
  \end{itemize}
\end{frame}


\subsection{Abstraction du matériel}
\begin{frame}
  \frametitle{Abstraction du matériel} 
Un système d'exploitation contient un ensemble de programmes dont le
rôle est de piloter les périphériques matériels. \pause 

Le système d'exploitation sert d'intermédiaire entre le
matériel et les autres applications pour :\pause
\begin{itemize}
  \item simplifier l'utilisation  du matériel\pause
  \item réglementer et protéger l'accès au matériel (intégrité)\pause
  \item donner corps aux abstractions utiles (le \alert{système de
    fichiers}, ou la notion d'utilisateur).\pause
\end{itemize}
Au tout début de la séquence de démarrage, le programme
\alert{d'amorçage} ou \emph{bootstrap}, un programme résidant dans un
mémoire spécifique, déclenche le chargement du système
d'exploitation en mémoire.
\end{frame}

\subsection{Gestionnaire de processus}
\begin{frame}
\frametitle{Gestionnaire de processus}

On appelle \alert{processus} un programme informatique en cours d'exécution. 

\pause
Le système d'exploitation assure la gestion des processus. Notamment:
\begin{itemize}
\item Le chargement des programmes à partir d'une mémoire de masse vers
  la mémoire de travail (et saut sur le point d'entrée du programme). \pause
\item Le partage du temps processeur. \pause
\item Le partage de la mémoire. \pause
\item Le système gère également les interruptions matérielles, la
  communication entre les processus etc.
\end{itemize}
\end{frame}

\begin{frame}
  \frametitle{Temps partagé, ordonnancement}
\pause
\begin{itemize}
\item  Sur un système multitâches, il y a un seul processeur (ou un petit
    nombre de processeurs) et de nombreux processus.
\pause
\item Chaque processus est exécuté tour à tour, pendant un intervalle de
    temps court, de l'ordre de 10 à 100 ms, après quoi il retourne en
    attente (jusqu'à ce que son tour revienne).  
\pause
 \item Ceci donne une illusion de simultanéité et de parallélisme.
\pause
\item Le temps partagé peut être \alert{coopératif} (c'est le programme du
  processus qui prévoit de rendre temporairement le processeur), ou
  \alert{préemptif} (le processus est interrompu par le système).
\pause
  \item  \alert{Algorithme d'ordonnancement}. C'est le système d'exploitation qui fixe l'ordre dans lequel ces
    processus sont exécutés. 
\pause
  \item Exemple d'ordonnancement : le tourniquet.
\pause
  \item Un peu de temps est perdu à la \alert{commutation de
      contexte}.
  \end{itemize}
\end{frame}


\begin{frame}
  \frametitle{Mémoire partagée, mémoire virtuelle}
  \begin{itemize}
\item La mémoire est en général virtualisée. \pause
\item Cela signifie que les applications n'ont pas d'accès direct à la
  mémoire physique. Les adresses mémoires manipulées par les processus
  ne sont pas directement les adresses physiques, elles sont
  virtuelles.\pause
  \item Un calcul d'adresse est effectué par le système pour passer
    des adresses virtuelles aux adresses physiques et
    réciproquement.\pause 
  \item Ceci permet de faire comme si chaque programme était seul
    résidant en mémoire et simplifie le travail de programmation.\pause
  \item Certaines portions inactives de mémoire de travail peuvent
    également être stockée sur une mémoire de masse (disque dur), de
    manière à étendre la quantité de mémoire de
    travail disponible (attention aux \alert{temps d'accès} !).
  \end{itemize}
\end{frame}

\subsection{Environnement de travail}
\begin{frame}
  \frametitle{Environnement de travail}
  Selon un point de vue extensif, le système est vu comme un ensemble
  de logiciels (navigateur web, éditeur de texte, etc.)  qui
  accompagnent le matériel et sont installés en même temps que le
  système d'exploitation proprement dit.\pause
\begin{itemize}
\item Environnement graphique (menus, bureau, corbeille) \pause
\item Services inter-applications (correction orthographique, annuaire, etc.)\pause
\item Multi-utilisateurs.\pause 
\end{itemize}

\structure{Remarque :} les mots de passe ne sont pas stockés, un chiffrement irréversible suffit.
\end{frame}


\section[Démos]{Démos et fin}
\begin{frame}
  \frametitle{Démos et fin}
\end{frame}
\end{document}
